\chapter{Эксперименты}\label{ch:5}

В работе проведено четыре последовательных эксперимента. Первый служил базовой проверкой работоспособности 2$\times$2 факторного плана из §2.2 на исходных гиперпараметрах (v1). Второй уже учитывал диагностические выводы первого и был выполнен с усовершенствованной конфигурацией (v2). Третий расширил линейку пятой архитектурой (AST v2). Четвёртый эксперимент задумывался как ещё одно усиление CNN, но оказался неудачным и в работу включён как академически ценный negative result. Сверх этого выполнен отдельный ablation по постобработке предсказаний сегментной головы (§5.8).

\section{Протокол оценки}\label{sec:5-1}

Качество модели оценивается стандартными для классификации метриками: accuracy, macro-F1 и поклассовые precision, recall, F1. Для структурной сегментации дополнительно считается Boundary F1 с допуском $\pm 3$ секунды. Это стандартная MIR-метрика, оценивающая, насколько правильно модель находит сами моменты переходов между секциями, безотносительно того, как она их маркирует.

Все три датасета разбиты на 70/15/15 (обучающую, валидационную и тестовую выборки) с зафиксированным сидом 42 (см. §3.3). Размеры тестовых сплитов после нарезки окон: примерно 11.6 тысяч окон у DEAM, 23.4 тысяч у Harmonix и 3.2 тысяч у GTZAN. Оценка делается один раз по чекпоинту с лучшей валидационной потерей, без подгонки порогов под тест.

В качестве точки отсчёта используется случайная базовая модель, выдающая равномерное распределение по классам. Для рассматриваемых задач это даёт 16{,}7\,\% по сегментной шестиклассовой задаче, 33{,}3\,\% для возбуждения и валентности и 10\,\% по жанру.

Все метрики приводятся по одному запуску с фиксированным сидом 42. Из-за стоимости полного обучения (около трёх часов на AST v2 и порядка часа на каждую из остальных конфигураций) в рамках данной работы не проводилась оценка статистической значимости различий между моделями: не вычислялись доверительные интервалы методом bootstrap по трекам, не проводилось повторных запусков с разными начальными сидами, не рассчитывалось стандартное отклонение метрик. В связи с этим небольшие различия в метриках (порядка 1--2~процентных пункта) следует интерпретировать с осторожностью, а полноценная количественная проверка значимости различий вынесена в направления дальнейшей работы (см. §7.3).

\section{Гиперпараметры обучения}\label{sec:5-2}

Базовая конфигурация v2 для всех архитектур, кроме AST, сведена в табл.~\ref{tab:hyperparams-v2}. Конфигурация AST отличается дифференциальными скоростями обучения и накоплением градиента и подробно описана в §4.5.

\begin{table}[H]
\centering
\caption{Гиперпараметры обучения (v2)}
\label{tab:hyperparams-v2}
\begin{tabular}{L{5cm}L{9cm}}
\toprule
Параметр & Значение \\
\midrule
Оптимизатор & AdamW, lr=1.0e-3, weight\_decay=5.0e-4 \\
Размер батча / эпохи / patience & 64 / 50 / 15 \\
Расписание скорости обучения & CosineAnnealingWarmRestarts (T\_0=10, T\_mult=2) \\
Функция потерь & Focal Loss, $\gamma=1.5$, max\_class\_weight=2.5 \\
MixUp $\alpha$ & 0.2 (per-window), 0.0 (BiLSTM) \\
Длина окна / шаг & 10.0~с / 1.0~с \\
SpecAugment & freq\_mask=27, time\_mask=60, n\_masks=2 \\
Dropout & 0.4 \\
Веса задач в общей потере & seg:1.5, aro\_cls:1.0, val\_cls:1.0, gen:1.0, *\_reg:0.5 \\
\bottomrule
\end{tabular}
\end{table}

\section{Эксперимент v1: исходный 2$\times$2 запуск}\label{sec:5-3}

Первая серия экспериментов прошла с базовой конфигурацией: четыре свёрточных блока (без пятого), 30 эпох обучения, расписание OneCycleLR, focal loss с $\gamma = 2.0$, верхний порог весов классов 4.0 и окно длиной 8 секунд. Результаты приведены в табл.~\ref{tab:v1-results}.

\begin{table}[H]
\centering
\caption{Результаты 2$\times$2 эксперимента в конфигурации v1}
\label{tab:v1-results}
\begin{tabular}{L{3.5cm}L{2.6cm}L{2.6cm}L{2.6cm}L{2.6cm}}
\toprule
Конфигурация & Segment acc/F1 & Arousal acc/F1 & Valence acc/F1 & Genre acc/F1 \\
\midrule
CNN per-window & 24.3 / 23.1 & 65.9 / 65.3 & 55.8 / 53.4 & 82.7 / 81.9 \\
CNN+BiLSTM & 25.9 / 26.8 & 53.0 / 51.8 & 49.5 / 48.4 & 45.5 / 42.5 \\
PANNs+Linear & 23.2 / 20.5 & 60.9 / 61.4 & 54.3 / 52.7 & 81.6 / 81.2 \\
PANNs+BiLSTM & 31.4 / 31.5 & 62.8 / 62.9 & 55.9 / 56.0 & 81.0 / 80.5 \\
\bottomrule
\end{tabular}
\end{table}

Первое же сравнение с базовой случайной линией показало, что все четыре конфигурации существенно выше случайного выбора, а самым слабым звеном везде оказалась структурная сегментация. Чтобы понять источник плохого качества, были рассмотрены поклассовые метрики CNN, на которой сегментная задача показала наихудшие результаты; распределение precision и recall по шести классам оказалось типичным для агрессивной компенсации дисбаланса. Самые частые классы Verse и Chorus имели высокую precision (порядка 70\%), но катастрофически низкий recall (10\% и 18\% соответственно). Редкие классы Intro, Bridge и Instrumental, наоборот, оказались сильно переопределены: модель ставила их везде, где сомневалась, что давало им recall выше 50\% при precision ниже 30\%. Концовка проседала по обеим осям (precision 4\%, recall 14\%). Поклассовое сравнение v1 и v2 на сегментной задаче приведено ниже в табл.~\ref{tab:v1-v2-per-class}. Для краткости отдельная таблица только для v1 опущена.

Главным виновником был выбран чрезмерно высокий верхний порог весов классов \code{max\_class\_weight=4.0}, давший редким классам слишком большое влияние. Отдельно бросилось в глаза катастрофическое падение жанра у CNN+BiLSTM до 45.5\% (явное переобучение на 699 трёхминутных треках GTZAN, которые после агрегации по последовательностям дают слишком мало обучающих примеров для BiLSTM).

Тепловая карта факторного дизайна (рис.~\ref{fig:factorial}) наглядно показывает, как два фактора плана (временное моделирование и качество признаков) проявились на v1.

\begin{figure}[H]
\centering
\includegraphics[width=0.62\textwidth]{fig_factorial_heatmap.png}
\caption{Тепловая карта 2$\times$2 факторного дизайна v1: временное моделирование по горизонтали, качество признаков по вертикали.}
\label{fig:factorial}
\end{figure}

\section{Эксперимент v2: улучшенный запуск}\label{sec:5-4}

На основе диагностики v1 в конфигурацию v2 была внесена серия изменений. Главные из них: добавлен пятый свёрточный блок с SE-модулем (для увеличения ёмкости), окно увеличено с 8 до 10 секунд, focal $\gamma$ снижен с 2.0 до 1.5, верхний порог весов классов уменьшен с 4.0 до 2.5, вес сегментной задачи в общей потере поднят с 1.0 до 1.5, обучение продлено до 50 эпох с patience 15, расписание скорости заменено с OneCycleLR на CosineAnnealingWarmRestarts, дополнительно введён временной джиттер $\pm 1.0$~с для устойчивости к границам сегментов. Запуск проводился единым прогоном, без отдельного анализа эффекта каждого изменения; общая ablation отдельных компонентов потери в работу не вошла из-за стоимости полного обучения. Полные результаты v2 приведены в табл.~\ref{tab:v2-results}.

\begin{table}[H]
\centering
\caption{Результаты 2$\times$2 эксперимента в конфигурации v2}
\label{tab:v2-results}
\begin{tabular}{L{3.5cm}L{2.6cm}L{2.6cm}L{2.6cm}L{2.6cm}}
\toprule
Конфигурация & Segment acc/F1 & Arousal acc/F1 & Valence acc/F1 & Genre acc/F1 \\
\midrule
CNN per-window & 35.3 / 29.5 & 62.0 / 61.0 & 51.7 / 51.1 & 83.9 / 83.3 \\
CNN+BiLSTM & 32.2 / 31.6 & 62.3 / 62.4 & 54.2 / 53.3 & 60.9 / 62.0 \\
PANNs+Linear & 24.0 / 21.7 & 61.2 / 61.6 & 53.8 / 53.3 & 81.9 / 81.4 \\
PANNs+BiLSTM & 32.9 / 32.7 & 64.9 / 64.0 & 55.0 / 53.0 & 77.7 / 77.9 \\
\bottomrule
\end{tabular}
\end{table}

Дельты относительно v1 распределились по архитектурам по-разному. CNN per-window получил +11.0/+6.4 п.п. на сегменте, заплатив за это $-3.9$ на возбуждении и $-4.1$ на валентности. CNN+BiLSTM прибавил +6.3 на сегменте, +9.3 на возбуждении и заметные +15.4 на жанре (то есть наконец оправился от переобучения v1). PANNs+Linear практически не сдвинулся, все дельты в пределах ±1.5 п.п. PANNs+BiLSTM прибавил +1.5 на сегменте и +2.1 на возбуждении, но потерял $-3.3$ на жанре.

Главное улучшение, ради которого и затевался переход к v2, действительно случилось: CNN прибавила 11 процентных пунктов по точности сегментной задачи, recall куплета вырос с 10 до 28\%, припева с 18 до 40\%. Конфигурация CNN+BiLSTM выровнялась и стала более конкурентной: рост на жанре +15 процентных пунктов означает, что переобучение на 699 треках GTZAN частично снято за счёт более жёсткого расписания обучения и более мягкой схемы взвешивания классов. Цена за приоритизацию сегмента просматривается тоже отчётливо: CNN потеряла около четырёх процентных пунктов на эмоциях, потому что её градиенты теперь сильнее тянуло в сторону структурной задачи.

Сводное поклассовое сравнение для CNN v1 против v2 на сегментной задаче приведено в табл.~\ref{tab:v1-v2-per-class}. Главный сдвиг здесь приходится на куплет и припев: распознавание самых частых и практически важных классов восстановилось.

\begin{table}[H]
\centering
\caption{Поклассовое сравнение сегментной задачи для CNN v1 и v2}
\label{tab:v1-v2-per-class}
\begin{tabular}{L{2.5cm}L{3cm}L{3cm}L{2cm}}
\toprule
Класс & v1 P/R/F1 & v2 P/R/F1 & $\Delta$ Recall \\
\midrule
Intro & 28.9 / 64.2 / 39.8 & 34.1 / 59.8 / 43.4 & $-4.4$ \\
Verse & 68.9 / 10.1 / 17.6 & 52.1 / 27.9 / 36.4 & +17.8 \\
Bridge & 19.2 / 51.5 / 28.0 & 17.3 / 39.8 / 24.1 & $-11.7$ \\
Chorus & 70.9 / 18.4 / 29.2 & 53.5 / 40.3 / 46.0 & +21.9 \\
Instrumental & 10.0 / 62.5 / 17.3 & 16.5 / 29.4 / 21.1 & $-33.1$ \\
Outro & 4.1 / 14.4 / 6.4 & 6.2 / 5.9 / 6.0 & $-8.5$ \\
\bottomrule
\end{tabular}
\end{table}

Результаты v2 позволяют дать предварительные ответы на первые четыре исследовательских вопроса. Одна модель действительно справляется со всеми четырьмя задачами одновременно: жанр у CNN v2 достиг 83.9\%, что укладывается в диапазон современных однозадачных моделей для GTZAN. Вклад BiLSTM на мультизадачной модели оказался неоднозначным: эмоции он действительно улучшает, но жанровую задачу у CNN+BiLSTM серьёзно подкашивает за счёт критически малой обучающей выборки в последовательностной формулировке. Замороженные PANNs не дают принципиального преимущества перед свёрткой, обученной с нуля, и это, в общем-то, ожидаемо: AudioSet ориентирован на разметку звуковых событий, а не музыкальной формы. И, наконец, ни временное моделирование, ни качество входных признаков по отдельности не оказались решающим фактором. Реальный прорыв принесла совершенно другая ось — переход на трансформерную архитектуру.

\section{Эксперимент AST v2: расширение трансформер-архитектурой}\label{sec:5-5}

После того как стало понятно, что в рамках 2$\times$2 факторного плана дальнейшие улучшения ограничены, эксперимент был расширен пятой архитектурой. AST~\cite{gong2021ast} был выбран потому, что одновременно даёт глобальное рецептивное поле (важно для структурной задачи) и богатую инициализацию весов от AudioSet (важно для общего качества). Конфигурация AST v2 описана в §4.5 и кратко повторяется тут: чекпоинт \code{MIT/ast-finetuned-audioset-10-10-0.4593}, около 87 миллионов параметров, из которых обучается порядка 14 миллионов (верхние 4 трансформерных слоя, LayerNorm, CLS, головы), дифференциальная скорость обучения и накопление градиента. Обучение заняло порядка трёх часов на 12-гигабайтной видеокарте.

Сводные метрики AST v2 на тестовых сплитах и сравнение с лучшей не-трансформерной моделью (CNN v2) приведены в табл.~\ref{tab:ast-v2-results}.

\begin{table}[H]
\centering
\caption{Результаты AST v2 на тестовых выборках}
\label{tab:ast-v2-results}
\begin{tabular}{L{3cm}L{2.5cm}L{2.5cm}L{4cm}}
\toprule
Задача & Accuracy & Macro-F1 & $\Delta$ к CNN v2 (по accuracy) \\
\midrule
Segment & 40.6 & 30.8 & +5.3 п.п. \\
Arousal & 60.6 & 60.4 & $-1.4$ п.п. \\
Valence & 51.1 & 50.4 & $-0.6$ п.п. \\
Genre & 81.5 & 80.8 & $-2.4$ п.п. \\
Среднее & 58.5 & 55.6 & +0.3 п.п. \\
\bottomrule
\end{tabular}
\end{table}

\begin{table}[H]
\centering
\caption{Поклассовые метрики сегментной задачи для AST v2}
\label{tab:ast-v2-per-class}
\begin{tabular}{L{3cm}L{2.5cm}L{2.5cm}L{2cm}}
\toprule
Класс & Precision & Recall & F1 \\
\midrule
Intro & 31.0 & 50.0 & 38.3 \\
Verse & 51.0 & 45.7 & 48.2 \\
Bridge & 20.1 & 22.9 & 21.4 \\
Chorus & 48.2 & 45.6 & 46.9 \\
Instrumental & 26.6 & 23.4 & 24.9 \\
Outro & 4.5 & 6.1 & 5.2 \\
\bottomrule
\end{tabular}
\end{table}

Ответ на пятый исследовательский вопрос получился содержательным. AST действительно превзошёл лучшую свёрточную модель именно там, где это было важнее всего: структурная сегментация улучшилась на 5{,}3~процентных пункта. Главный источник прироста виден из поклассовых метрик AST v2 (табл.~\ref{tab:ast-v2-per-class}). Recall на куплете и припеве у AST сбалансированы заметно лучше: 45{,}7\,\% против 27{,}9\,\% у CNN на куплете, 45{,}6\,\% против 40{,}3\,\% на припеве. Глобальное рецептивное поле помогает модели увидеть всю композицию целиком, а предобучение на AudioSet даёт стартовое представление о низкоуровневых акустических признаках, которое сложно выучить с нуля на 522 структурных треках.

Качество остаётся существенно неравномерным по классам. AST v2 уверенно различает крупные и частые секции (Verse F1 = 48{,}2, Chorus F1 = 46{,}9), но плохо справляется с редкими и акустически размытыми классами: Bridge F1 = 21{,}4, Instrumental F1 = 24{,}9, Outro F1 = 5{,}2. Концовка по-прежнему распознаётся хуже всего и фактически на уровне случайной модели, что отражает сочетание малой представленности класса в обучающих данных и его акустической близости к Verse и Instrumental. На задаче жанра AST уступает CNN v2 на 2.4 процентных пункта. Скорее всего, для десятиклассовой классификации на коротких 30-секундных фрагментах GTZAN inductive bias свёртки оказался эффективнее. С точки зрения практической развёртки за это улучшение приходится заплатить: AST занимает в памяти около 574 МБ против 21 МБ у CNN, а инференс на одном треке требует примерно в четыре раза больше времени.

\section{Эксперимент CNN v3: неудачное усиление}\label{sec:5-6}

После того как стало понятно, что AST v2 выигрывает на сегментной задаче за счёт глобального контекста, была проверена гипотеза: возможно, CNN v2 также можно немного улучшить на сегментной задаче за счёт тройного усиления компенсаций дисбаланса. Вес головы сегментации был повышен с 1.5 до 3.5, показатель focal loss с 1.5 до 2.5, верхний порог весов классов с 2.5 до 3.5. Результаты получились прямо противоположные ожиданиям: на сегменте accuracy обвалилась на 13.2 процентных пункта (с 35.3 до 22.1), жанр упал на 4.8 пункта (с 83.9 до 79.1), и только эмоциональные задачи получили скромный прирост +2.5/+2.4. Среднее по четырём задачам ушло вниз на 3.3 процентных пункта (полные числа CNN v3 включены в сводную табл.~\ref{tab:final-comparison} строкой 7).

Разбор поведения модели показал быстрое переобучение: лучший валидационный лосс достигнут уже на четвёртой эпохе, early stopping сработал на девятнадцатой. После применения верхнего порога 3.5 веса классов сегментной задачи приняли вид [Intro:3.05, Verse:0.44, Chorus:0.62, Bridge:3.5, Instrumental:3.5, Outro:3.5], то есть редкие классы получили в шесть-восемь раз больший вклад в потерю, чем доминирующие. В результате модель сместилась в сторону чрезмерного предсказания редких классов настолько, что метрики по самым частым (Verse, Chorus) ухудшились вместе с общей точностью примерно на 13~процентных пунктов.

Этот эксперимент включён в работу сознательно как аккуратно поставленный отрицательный результат (negative result), демонстрирующий следующее: для свёрточной архитектуры рассматриваемых размеров на используемых данных конфигурация v2 уже находится близко к локальному оптимуму компенсаций дисбаланса, и дальнейшее усиление приводит к ухудшению. Принципиальное улучшение сегментной задачи на этом этапе оказывается достижимо не за счёт тонкой настройки потери, а за счёт смены архитектуры (что и сделано в AST v2).

\section{Сводное сравнение всех моделей}\label{sec:5-7}

Итоговое сравнение всех семи экспериментальных конфигураций приведено в табл.~\ref{tab:final-comparison}. Колонка «Параметров» для предобученных моделей (3--6) указывает только обучаемую часть.

\begin{table}[H]
\centering
\caption{Финальное сравнение всех конфигураций по тестовым метрикам, accuracy в \%}
\label{tab:final-comparison}
\footnotesize
\setlength{\tabcolsep}{4pt}
\begin{tabular}{L{0.5cm}L{3cm}L{2cm}L{1.5cm}L{1.5cm}L{1.5cm}L{1.5cm}L{1.5cm}}
\toprule
\# & Модель & Параметров & Segment & Arousal & Valence & Genre & Среднее \\
\midrule
1 & CNN v1 & $\sim$400K & 24.3 & 65.9 & 55.8 & 82.7 & 57.2 \\
2 & CNN v2 & $\sim$800K & 35.3 & 62.0 & 51.7 & 83.9 & 58.2 \\
3 & CNN+BiLSTM v2 & $\sim$1.2M & 32.2 & 62.3 & 54.2 & 60.9 & 52.4 \\
4 & PANNs+Linear v2 & $\sim$210K & 24.0 & 61.2 & 53.8 & 81.9 & 55.2 \\
5 & PANNs+BiLSTM v2 & $\sim$750K & 32.9 & 64.9 & 55.0 & 77.7 & 57.6 \\
6 & AST v2 & $\sim$14M & 40.6 & 60.6 & 51.1 & 81.5 & 58.5 \\
7 & CNN v3 (failed) & $\sim$800K & 22.1 & 64.5 & 54.1 & 79.1 & 54.9 \\
\bottomrule
\end{tabular}
\end{table}

Победители по задачам распределились следующим образом. На структурной сегментации лидирует AST v2 с 40.6\% точности. По эмоциональным задачам максимальные числа показывает PANNs+BiLSTM (arousal 64.9\%, valence 55.0\%). Это совпадает с известным наблюдением, что для предсказания эмоции хорошо работает сочетание сильных аудио-признаков AudioSet с временным контекстом всей композиции. По жанру лидирует CNN v2 (83.9\%). На коротких 30-секундных фрагментах GTZAN inductive bias свёртки оказывается эффективнее самовнимания. По среднему результату по всем четырём задачам AST v2 опережает CNN v2 минимально, на 0{,}3~процентных пункта (58{,}5\,\% против 58{,}2\,\%), что находится в пределах статистической погрешности. Выбор AST v2 как финальной модели прототипа оправдан прежде всего наибольшим качеством по структурной сегментации, которая является приоритетной задачей системы; для сценариев с приоритетом жанровой или эмоциональной классификации более уместными выбраны бы CNN v2 или PANNs+BiLSTM соответственно.

\begin{figure}[H]
\centering
\includegraphics[width=0.65\textwidth]{fig_comparison_5configs.png}
\caption{Сравнение пяти архитектур по четырём задачам (v2 + AST v2).}
\label{fig:comparison-5}
\end{figure}

\section{Ablation: вклад постобработки сегментной головы}\label{sec:5-8}

Для оценки вклада постобработки выполнен отдельный ablation на тестовом сплите Harmonix Set с предсказаниями AST v2 (табл.~\ref{tab:ablation-postprocess}). Каждая строка добавляет один механизм поверх предыдущей.

\begin{table}[H]
\centering
\caption{Ablation постобработки на тестовой выборке Harmonix Set}
\label{tab:ablation-postprocess}
\begin{tabular}{L{5cm}L{3.5cm}L{2.5cm}L{2.5cm}}
\toprule
Метод & Сегментов (в среднем) & Boundary F1 $\pm$3~с & Section F1 \\
\midrule
Argmax (без постобработки) & 28.4 & 0.31 & 0.32 \\
+ Median filter (kernel=5) & 14.2 & 0.36 & 0.34 \\
+ Viterbi ($\lambda = 1.0$) & 9.8 & 0.42 & 0.37 \\
+ Position priors (penalty=10) & 8.7 & 0.43 & 0.38 \\
+ Min-duration merge (6~с) & 7.5 & 0.45 & 0.39 \\
\bottomrule
\end{tabular}
\end{table}

Видно, что самый большой вклад даёт переход от медианного фильтра к декодеру Витерби, что добавляет сразу 6 процентных пунктов к Boundary F1. Это и понятно: первичное сглаживание убирает только локальный шум, а Витерби уже учитывает глобальную последовательность меток с штрафом за переходы. Полный пайплайн (Витерби с позиционными приорами и слиянием коротких сегментов по минимальной длительности) даёт суммарный прирост 14 процентных пунктов к Boundary F1 и 7 процентных пунктов к Section F1 относительно сырого argmax. Среднее число сегментов в трёхминутном треке снижается с 28 до 7--8, что уже укладывается в диапазон реальной разметки в Harmonix Set (среднее по корпусу 8--12 сегментов).
